% Source: https://artofproblemsolving.com/community/c3h2487867_square_root_inside_a_square_root
% Stable identity: aops-topic-2487867-q01
\textbf{(a)} Since $Q\ge0$ we have $R\le P$, and since $Q\le P^2$ the number $R$ is real and $R\ge0$. Hence $P+R\ge P-R\ge0$, so both square roots on the right are defined and the right-hand side is nonnegative. Also $\sqrt Q\le P$, so the left-hand side is defined and nonnegative.

Squaring the right-hand side gives
\[\dfrac{P+R}{2}+\dfrac{P-R}{2}-2\sqrt{\dfrac{(P+R)(P-R)}{4}}=P-\sqrt{P^2-R^2}=P-\sqrt Q,\]
which is the square of the left-hand side. Two nonnegative numbers with equal squares are equal, which proves the identity. $\blacksquare$

\textbf{(b)} Move the factor inside: $x=\sqrt{8-2\sqrt7}=\sqrt{8-\sqrt{28}}$. Here $P=8$, $Q=28$ and $R=\sqrt{64-28}=6$, so
\[x=\sqrt{\dfrac{14}{2}}-\sqrt{\dfrac{2}{2}}=\boxed{\sqrt7-1}.\]
Check: $\left(\sqrt7-1\right)^2=8-2\sqrt7$ and $\sqrt7-1>0$.

\textbf{(c)} Similarly $y=\sqrt{15-3\sqrt{21}}=\sqrt{15-\sqrt{189}}$, with $P=15$, $Q=189$ and $R=\sqrt{225-189}=6$. So
\[y=\sqrt{\dfrac{21}{2}}-\sqrt{\dfrac{9}{2}}=\dfrac{\sqrt{42}-\sqrt{18}}{2}=\boxed{\dfrac{\sqrt{42}-\sqrt{18}}{2}}=\dfrac{\sqrt{42}-3\sqrt2}{2}.\]

\emph{Remark.} In general $a\left(b-\sqrt c\right)=ab-\sqrt{a^2c}$, so $P=ab$ and $Q=a^2c$. The formula always holds, but it removes the nested radical only when $P^2-Q=a^2\left(b^2-c\right)$ is a perfect square; otherwise $R$ is itself a square root and the expression is not simpler.
